%% ---------------------------------------------------------------------------
%% intro.tex
%%
%% Introduction
%%
%% $Id: intro.tex 1477 2010-07-28 21:34:43Z palvarado $
%% ---------------------------------------------------------------------------

\chapter{Introduction}
\label{chp:intro}

We are said to be living in a \textit{New Space} era, where access to space is open to more countries, companies, universities, and people than ever.
This democratization of space has been driven by progressively decreasing barriers to entry for space systems development and space launches in terms of time horizon, technological feasibility, and economic cost.
The factors at play include reusable launch platforms, standardization of space system and payload architectures, the integration of commercial-off-the-shelf (COTS) components, and the adoption of more efficient and effective design and development approaches.

The CubeSat format within the nanosatellite class is a prime example of architecture standardization.
These systems consist of modular units of 10-by-10-by-10 centimeter cubes at around 1.33 kilograms, which have been proven to be integratable at up to 16 units \cite{kulu_nanosats_database} and have supported a plethora of scientific and commercial payloads ranging from orbital laboratories to internet network nodes.

Unfortunately, CubeSats historically suffer from a high failure rate, and thus not only fall short of fulfilling mission needs but become a worrying source of space debris at the crowded low earth orbits that they usually launch to.
The causes can be traced back to a lack of expertise, coordination, or consistency in the engineering teams behind the mission, often based at universities that are new entrants to the sector as opposed to incumbents with space heritage, or, at a technical level, attributed to electrical power systems (EPS) as a centralized point of unrecoverable failure.

This project takes place at the (ASCEND) research group at the École de technologie supérieure in Montréal, Canada.
Dr. Jesús González Llorente, ASCEND lead an academic advisor of this project, pursues a line of research concerned with reliable space systems, and has authored multiple academic publications targeting the design and validation of nanosatellite electrical power systems.

ASCEND is pursuing participation in a CubeSat Initiative in Canada for STEM (CUBICS) mission co-developed across two universities and involving mostly student volunteers in engineering and coordination roles.
This endeavor requires a wide and deep exploration of solution candidates, coupled with robust tradeoff analysis, verification, and validation of said candidates to deliver a reliable CubeSat.
Furthermore, the CUBICS financial and time budget requires that any requirement non-conformities need to be detected as early as possible, preferrably prior to hardware fabrication which is one of the costliest phases of the entire engineering lifecycle, and is associated with a large iteration period.
All of the above needs to achieved by a highly motivated and capable yet inexperienced cohort of student volunteers.

The current project is intended as a precedent for the CUBICS ambition, which is referred to as capability development. It covers the electrical power subsystem in its lifecycle phases prior to commissiong or developing hardware, and is founded in the experience gained in previous EPS design efforts at ASCEND.
The problem that motivated the current project is:

\begin{quote}
Evidence fragmentation during the early lifecycle phases of nanosatellite electrical power systems hinders the traceability between needs and requirements, architecture decisions, and the analytical evidence that supports them. Furthermore, it hinfers the reproducibility of analyses and the capability to detect design inconsistencies before advacing to physical verification and validation phases.
\end{quote}

This project applied the Object-Oriented Systems Engineering Method to the above problem.
The as-is EPS design and development practices of the stakeholders and their shortcomings were assessed, and a to-be set of capabilities was derived from best practices identified in the academic literature.
After multiple discussions with stakeholders, a set of concerns about the as-is capabilities and the potential transition to the to-be capabilities was constructed.
Identified concerns relate to the high effort to define verification activities for EPS component requirements, the high effort to maintain traceability from requirements to design decisions, the high cost and effort of redesign after hardware fabrication, and several other factors.

The main objective of this project is:

\begin{quote}
To develop a model-based systems engineering framework oriented at the early justification of nanosatellite electrical power systems through traceable and reproducible evidence.
\end{quote}

An MBSE framework is defined following Buruso et al. \cite{buruso2026modelbased} as the integration of MBSE methods, MBSE languages, and tools within and throughout the relevant lifecycle phases.
The prescription of MBSE in the main objective is a constraint imposed by ASCEND expertise and capability development interests.
Therefore, multiple methods, languages, and tools from within the ample MBSE literature were considered as candidate building blocks for the framework.

The secondary objectives of this project are: (1) to characterize the needs, requirements, interfaces and technical criteria of nanosatellite EPS; (2) to define the operational flow of modeling, analysis, and evidence management of the MBSE framework; (3) to model a generic EPS architecture with functions, logical and physical components, interfaces, modes and states, and requirements; (4) to characterize early EPS verification and validation within the EPS framework; (5) to evaluate the framework on two contrasting EPS case studies.
Therefore, the current project does not aim to only design or develop an EPS for ASCEND, but to improve the EPS design and development capabilities at the research group, so that they can be leveraged to bring the future CUBICS mission to fruition.

The framework will be verified against requirements such as retaining forward and reverse links between requirements and component allocations, enabling the reconstruction of mission analyses, and identifying change impacts.
Verification will combine catalogue inspection, two mission case studies, and two individual user trials.
The BIRDS-4 and OreSat case studies will use published power budgets and component data to check EPS mass and volume estimates, power and energy budgets, sizing, component and architecture alternatives, and electrical limits against independently prepared references.
The user trials will cover requirement and allocation traceability, analysis records and reproduction, change impact assessment, engineering guidance, revision control, COTS recommendations, and mission input management.
An ASCEND PhD candidate and a student club undergraduate will each perform 11 tasks, with each task required to take less than $60/11$~min per participant.

The remainder of this document is structured as follows: Chapter~\ref{ch:marco} presents the theoretical background, with particular attention to normative documents and state-of-the-art MBSE research. Chapter~\ref{ch:solucion} describes the OOSEM analysis of stakeholder needs and their formalization as verifiable framework requirements that can be fulfilled by multiple architecture alternatives. Chapter~\ref{ch:design-proposal} details the design of the chosen alternative, in terms of its CubeSat and EPS models, orbit characterization, and power budget, component selection, and verification method implementations. Chapter~\ref{ch:results} covers results and analysis, obtained via two case studies as well as user trials. Chapter~\ref{ch:conclusions} states the conclusions and future work.

% The intro must close with a paragraph stating what was the main engineering contribution of the project. This is best written when the project is complete, so not yet.
The main engineering contribution of this project is the publication of the full codebase \cite{repo} of an MBSE framework that can mitigate evidence fragmentation during the early lifecycle phases of nanosatellite electrical power systems owing to it enabling traceable and reproducible evidence.
The framework imposes no operation costs as it is implemented using Python and Github, and has been verified against BIRDS-4 and OreSat with regard to the correctness of its analyses.
User trials showed that the stakeholders became more agile when managing EPS design evidence when using the framework, and it has potential to be extended to other nanosatellite subsystems.
